\documentclass[11pt,a4paper]{article}
\usepackage[UTF8]{ctex}
\usepackage{geometry}
\usepackage{enumitem}
\usepackage{titlesec}
\usepackage{xcolor}
\usepackage{hyperref}

\geometry{left=2.5cm,right=2.5cm,top=2.5cm,bottom=2.5cm}

% 定义标题格式
\titleformat{\section}{\large\bfseries\color{blue!70!black}}{\thesection}{1em}{}
\titleformat{\subsection}{\normalsize\bfseries\color{blue!50!black}}{\thesubsection}{1em}{}

\begin{document}

\begin{center}
    {\LARGE\bfseries 项目经历：基于 NURBS 的球壳多面体铺砖参数化建模系统}\\[0.8em]
    {\large SphereShellTile：Sphere Shell Tiling with NURBS Patches}\\[1.2em]
    \textbf{技术栈：} MATLAB \quad|\quad NURBS Toolbox \quad|\quad 计算几何 \quad|\quad 等几何分析 (IGA)
\end{center}

\section{项目背景与目标}
在等几何分析（Isogeometric Analysis, IGA）与高精度 CAD/CAE 集成的研究中，球面及球壳结构的精确参数化表示是基础难题。传统上，球面无法被单一 NURBS 曲面片无畸变地覆盖，必须采用\textbf{多片拼接（multipatch）}策略。本项目旨在构建一套完整的 MATLAB 工具链，实现从多面体对称性出发的球面/球壳 NURBS 铺砖（tiling）参数化，并输出标准化的几何数据，供后续数值求解器直接调用。

\section{核心贡献}

\subsection{算法实现：从顶点到 NURBS 曲面片}
基于 Dedoncker 等人的理论工作 \textit{``Bézier tilings of the sphere and applications in benchmarking multipatch isogeometric methods''}，独立实现了核心函数 \texttt{MakeTile.m}：
\begin{itemize}[leftmargin=2em]
    \item 接收球面上任意 4 个有序顶点（凸四边形），通过\textbf{球极平面投影}（Stereographic Projection，参考 Cobb, 1988）映射至 $z=-r$ 平面；
    \item 在投影平面上构造以\textbf{大圆弧}为边界的 $5\times5$ NURBS 控制网，利用配置点（collocation）方法求解线性系统，反投影得到球面上的双四次（bi-quartic）NURBS 曲面片；
    \item 支持\textbf{符号解}与\textbf{数值解}两种模式，兼顾精度与计算效率。
\end{itemize}

\subsection{多面体对称扩展：覆盖多种分辨率需求}
基于立方体单片的参数化结果，通过旋转矩阵自动化组装，系统性地支持了 6 种多面体对称铺砖方案，覆盖从低分辨率到高分辨率的完整谱系：
\begin{center}
\begin{tabular}{lcc}
\hline
\textbf{多面体类型} & \textbf{面片数量} & \textbf{主要文件} \\
\hline
Cube（立方体） & 6 & \texttt{Multi\_SphereFromCube.m} \\
Rhombic dodecahedron（菱形十二面体） & 12 & \texttt{Multi\_SphereFromRDode.m} \\
Icosahedron（二十面体） & 20 & \texttt{Multi\_SphereFromIco.m} \\
Deltoidal icositetrahedron（鸢形二十四面体） & 24 & \texttt{Multi\_SphereFromDIco.m} \\
Rhombic triacontahedron（菱形三十面体） & 30 & \texttt{Multi\_SphereFromRTria.m} \\
Deltoidal hexecontahedron（鸢形六十面体） & 60 & \texttt{Multi\_SphereFromDHexe.m} \\
\hline
\end{tabular}
\end{center}
每种多面体均配有预计算的控制点数据（\texttt{.mat}）与独立的加载函数（\texttt{*Load.m}），确保结果可复现、可扩展。

\subsection{球壳体参数化：从曲面到三维 NURBS 体}
区别于仅参数化球面的常规做法，本项目进一步实现了\textbf{球壳三维体（NURBS volume）}的构造：
\begin{itemize}[leftmargin=2em]
    \item 通过设定内外两层半径（如 $r_{\text{inner}}=98$, $r_{\text{outer}}=100$），分别生成内外 NURBS 曲面片；
    \item 沿径向（第三参数方向 $u$）将两层控制点拼接，构造次数为 $(4,4,1)$ 的 NURBS 体块（volume patch），完美描述具有一定厚度的球壳几何；
    \item 开发了完整的球壳专用驱动脚本 \texttt{driver\_shell.m}（基于立方体 6 片）与 \texttt{driver\_shell\_Hemisphere.m}（基于半球 2 片），可直接生成并可视化球壳体控制网格。
\end{itemize}

\subsection{数据交换接口：标准化 YAML 几何输出}
为了打通 MATLAB 原型与外部工程求解器之间的数据链路，编写了 \texttt{writeYAML.m} 输出模块：
\begin{itemize}[leftmargin=2em]
    \item 自动提取每个 NURBS 体块的节点向量（\texttt{GLOBAL\_S/T/U}）、多项式次数（\texttt{DEGREE\_S/T/U}）以及所有控制点坐标与权重；
    \item 按照 $s \to t \to u$ 的规范顺序展平控制点数组，生成结构化的 \texttt{patch*.yml} 文件；
    \item 输出格式严格、可读性强，可直接被基于 YAML 配置的 IGA 或有限元框架解析导入。
\end{itemize}

\section{技术亮点与创新}
\begin{enumerate}[leftmargin=2em]
    \item \textbf{理论到代码的完整落地：} 将高阶球面 Bézier tiling 的数学理论转化为工程可用的 MATLAB 代码，填补了文献中缺乏开源实现工具的空白。
    \item \textbf{统一的多面体铺砖框架：} 以旋转对称性为核心，构建了可复用的 multipatch 组装模式，使不同面片数量的球面参数化可在同一架构下快速扩展。
    \item \textbf{壳体结构的三维 NURBS 参数化：} 不仅停留在曲面层面，而是将厚度维度纳入参数化，生成真正的 NURBS 体，为壳体等几何分析提供了精确的几何前处理基础。
    \item \textbf{可视化与数据输出一体化：} 每个驱动脚本均包含控制点网格与曲面/体的三维可视化，配合标准化 YAML 导出，实现了“建模—可视化—数据交付”的完整闭环。
\end{enumerate}

\section{成果与影响}
\begin{itemize}[leftmargin=2em]
    \item 成功生成了从 2 片（半球壳）到 60 片（鸢形六十面体）多种精度的球壳 NURBS 几何模型；
    \item 输出的 YAML 文件已被用于后续等几何分析求解器的基准测试（benchmarking）与多片拼接算法验证；
    \item 代码结构模块化、注释完整，为团队后续开展复杂曲面/multipatch IGA 研究提供了可直接复用的基础平台。
\end{itemize}

\end{document}
